% ----------------------
\subsection{Section 121 (20 March 1839)}
% ----------------------
	\label{DC121}
	
	% -----
	\subsubsection{The original sources, and table of verse locations in the original sources}
	% -----
	
		This section of the DC is a composite of passages drawn from two different letters written in late March 1839, from a jail in Missouri. The second letter is a continuation of the first. Orson Pratt selected the passages from the letters to be added to the DC in 1876, but he didn't choose everything. The following table shows the verses and then the passages in the original letters, based on my divisions in the margin.
		
		\begin{table}[H]
			\centering
			\begin{tabular}{p{4cm}p{10cm}}
				\textbf{Verses in DC 121} & \textbf{Original source} \\
				1--33 & \hyperref[EMS-1839-JS-20-MAR-1839]{Letter to the Church and Edward Partridge, 20 March 1839} \\
				34--46 & \hyperref[EMS-1839-JS-CIR-22-MAR-1839]{Letter to Edward Partridge and the Church, circa 22 March 1839} \\
				\textbf{Verses} & \textbf{Locations in original letters} \\
				1--6 & \hyperlink{EMS-1839-JS-20-MAR-1839-f}{f--h} (first letter) \\
				7--25 & \hyperlink{EMS-1839-JS-20-MAR-1839-o}{o--t} (first letter) \\
				26--32 & \hyperlink{EMS-1839-JS-20-MAR-1839-bb}{bb--dd} (first letter) \\
				33 & \hyperlink{EMS-1839-JS-20-MAR-1839-ee}{ee} (first letter) \\
				34--46 & \hyperlink{EMS-1839-JS-CIR-22-MAR-1839-d}{d--h} (second letter) \\
			\end{tabular}
		\end{table}
	
	% -----
	\subsubsection{Historical Background}
	% -----
		\label{DC121-HB}
		
		% --
		\paragraph{JSP}
		% --
	
			``In the Clay County, Missouri, jail on 20 March 1839, JS dictated a letter addressed to Bishop Edward Partridge; church members in Quincy, Illinois; and the Saints `scattered abroad.' The letter was the second general epistle JS directed to the church while in the jail, with the first missive composed on 16 December 1838. As the main body of Latter-day Saints was relocating from Missouri to Illinois and Iowa Territory, JS apparently envisioned writing a series of general epistles in March 1839 to offer guidance and instruction in the wake of the catastrophic changes of the previous year. The immediate catalyst for the 20 March letter was the arrival of Latter-day Saint David W. Rogers the previous evening. Rogers brought the prisoners a packet that contained letters from their families and friends, a letter from Illinois land speculator Isaac Galland, and `the documents and papers sent by the authorities at Quincy.'
			
			``Stylistically, the 20 March letter is reminiscent of the apostle Paul's epistles in the New Testament. Paul frequently named his companions in opening greetings and utilized the first-person plural voice even though he was the primary author of the letters. In a similar fashion, the 20 March 1839 letter opens with greetings from JS `in company with his fellow prisoners'; the body of the letter consistently employs the first-person plural---`we,' `our,' and `us'---with the exception of one portion presented in the voice of Deity; and all the prisoners signed the letter. JS was the principal author,%
				\footnote{%
					% \label{}%
					According to the JSP, we know this because of what JS says in his \hyperref[EMS-1839-JS-21-MAR-1839]{letter} to Emma the next day.
					}
			although conversations with the other prisoners may have contributed to the letter's ideas and themes. It is unknown who acted as scribe for the dictation draft, which is apparently not extant. The version featured here, which contains errors usually associated with copying, was inscribed by Alexander McRae and Caleb Baldwin. After McRae finished copying the last portion of the letter, JS and the other men signed the copy. At some point, JS made minor corrections and additions.
			
			``Following the opening greeting, the epistle contains an extended meditation on the Latter-day Saints' recent sufferings and the prisoners' frustrations in jail. This part of the missive includes a prayer in which JS pleads with God to deliver the Saints from their oppressors. The subject of the letter appears to shift with the acknowledgment of receiving letters from Edward Partridge, Don Carlos Smith, and Emma Smith, but this narrative actually continues the meditation on the meaning of persecution, revealing that reading the letters dissolved feelings of bitterness and opened JS's heart to receive inspiration. Then, the voice of the letter transitions from that of the prisoners to that of the Lord providing an answer to the letter's earlier prayer, explaining the deeper significance of the Saints' persecutions and pronouncing judgments against the church's enemies.
			
			``The second part of the letter addresses challenges the church faced in moving forward, such as deciding where the Saints should settle. JS declined to either approve or reject Isaac Galland's offer to sell land to the church; instead, JS said that church leaders in Quincy should make that decision in future conferences and should forward minutes of the proceedings to JS for approval. The letter also advises the Quincy church leaders to eschew `an aspiring spirit' that had previously prevailed over `milder councils,' causing much suffering and death among the Saints. Additionally, the epistle contains counsel on how to seek revelation and guidance; this counsel is followed by strong affirmation that persecution would not hinder the work of God. Like other missives JS composed in the Clay County jail, this letter incorporates multiple biblical allusions. Near the close, the letter signals that another general epistle was forthcoming.
			
			``Although the letter's greeting is directed to the church in general and to Partridge in particular, JS sent the missive to his wife Emma because he wanted her `to have the first reading of it.' In a letter he wrote to her the following day, he informed her, `I have sent an Epistle to the church,' suggesting the epistle had already left the jail. The letter was probably carried from the jail by a church member, perhaps Alanson Ripley, who Hyrum Smith noted was visiting the jail on 20 March and was `going to start back this after noon' to Far West, Missouri. Ripley told the prisoners that he could send their letters to Illinois `a mediately by some of the brethren.' It remains unclear who transported the missive from Missouri to Illinois or when it arrived in Quincy. On 10 April 1839, Sidney Rigdon and Ripley wrote separate letters to JS and the other prisoners; both messages contain possible allusions to the general epistle, suggesting church members had received and read the epistle by that date. On 11 April, Mary Fielding Smith wrote to her husband, Hyrum Smith, stating she had read the epistle and that it was `food to the hungrey.' The 20 March epistle circulated widely among the Latter-day Saints in the months after its arrival in Illinois, as indicated by the extant copies in the handwriting of Partridge and Albert Perry Rockwood. In addition, the \textit{Times and Seasons} published an edited version in 1840,%
				\footnote{%
					% \label{}%
					It was first published in the 1876 DC, after Orson Pratt extracted some passages. BY was involved in the production of that edition of the DC but, according to HDDC 1566--1567, it's not clear by what criteria he selected them (and I think it also isn't clear what amount of oversight BY had for the entire thing).
					}
			extending the letter's circulation to the Saints `scattered abroad.'''
			
		% --
		\paragraph{HDDC} 
		% --
		
			HDDC 1566--1574 has some more background, including extracts of newspaper editorials from different sources that discuss the publication history of certain passages of this letter in various places. It's a lot of useless jawing now.
			
	% -----
	\subsubsection{Versions}
	% -----
		\label{DC121-V}
		
		See the \href{https://drive.google.com/file/d/1goAvNz8jS4R8slYfMG_db55Ty2z_vmgK/view?usp=sharing}{sections comparison} document.
	
		\begin{compactdesc}
			\item[JSP] \hyperref[EMS-1839-JS-20-MAR-1839]{20 March 1839}
			\item[BC] ---
			\item[DC 1835] ---
			\item[DC 1844] ---
			\item[TS] 1:7, 9 (May and July 1840), 99--104 and 131--134
			\item[MS] 1:8 (December 1840), 193--197
			\item[MS] 5:5 (October 1844), 69--72
			\item[DN] 4:4, 5 (26 January and 2 February 1854), 13 and 17
			\item[MS] 17:4, 6 (27 January and 10 February 1855), 52 and 84--88
		\end{compactdesc}
	
	% -----
	\subsubsection{Section 121:1} % *DC121-1 
	% -----
		\label{DC121-1}

			\footnote{%
				% \label{}%
				This section divides into five total chunks which come from the two original letters. You can see the table above which shows all the chunks, but the first one begins here, and continues until the end of verse 6. In the original that's \hyperlink{EMS-1839-JS-20-MAR-1839-f}{f--h}.
				}% 
		O GOD, where art thou?  And where is the pavilion that covereth thy hiding place?

		% \begin{framed}
		% \scriptsize
			% \begin{compactdesc}
				% \item[JSP] 
				% % \item[BC] 
				% % \item[]
			% \end{compactdesc}
		% \end{framed}

	% -----
	\subsubsection{Section 121:2} % *DC121-2 
	% -----
		\label{DC121-2}

		How long shall thy hand be stayed, and thine eye, yea thy pure eye,%
			\footnote{%
				% \label{}%
				Reminds me of \hyperref[DC131-8]{DC 131:8}.
				}
		behold from the eternal heavens the wrongs of thy people and of thy servants, and thine ear be penetrated with their cries?

		% \begin{framed}
		% \scriptsize
			% \begin{compactdesc}
				% \item[JSP] 
				% % \item[BC] 
				% % \item[]
			% \end{compactdesc}
		% \end{framed}

	% -----
	\subsubsection{Section 121:3} % *DC121-3 
	% -----
		\label{DC121-3}

		Yea, O Lord, how long shall they suffer these wrongs and unlawful oppressions, before thine heart shall be softened toward them, and thy bowels be moved with compassion toward them?

		% \begin{framed}
		% \scriptsize
			% \begin{compactdesc}
				% \item[JSP] 
				% % \item[BC] 
				% % \item[]
			% \end{compactdesc}
		% \end{framed}

	% -----
	\subsubsection{Section 121:4} % *DC121-4 
	% -----
		\label{DC121-4}

		O Lord God Almighty, maker of heaven, earth, and seas, and of all things that in them are, and who controllest and subjectest the devil, and the dark and benighted dominion of Sheol%
			\footnote{%
				% \label{}%
				This is the only use of ``Sheol'' in all the scriptures. The Hebrew word is \he{+s:'Ol}, HALOT 1368b.
				}%
		---stretch forth thy hand; let thine eye pierce; let thy pavilion be taken up; let thy hiding place no longer be covered; let thine ear be inclined; let thine heart be softened, and thy bowels moved with compassion toward us.

		% \begin{framed}
		% \scriptsize
			% \begin{compactdesc}
				% \item[JSP] 
				% % \item[BC] 
				% % \item[]
			% \end{compactdesc}
		% \end{framed}

	% -----
	\subsubsection{Section 121:5} % *DC121-5 
	% -----
		\label{DC121-5}

		Let thine anger be kindled against our enemies; and, in the fury of thine heart, with thy sword avenge us of our wrongs.

		% \begin{framed}
		% \scriptsize
			% \begin{compactdesc}
				% \item[JSP] 
				% % \item[BC] 
				% % \item[]
			% \end{compactdesc}
		% \end{framed}

	% -----
	\subsubsection{Section 121:6} % *DC121-6 
	% -----
		\label{DC121-6}

		Remember thy suffering saints, O our God; and thy servants will rejoice in thy name forever.

		% \begin{framed}
		% \scriptsize
			% \begin{compactdesc}
				% \item[JSP] 
				% % \item[BC] 
				% % \item[]
			% \end{compactdesc}
		% \end{framed}

	% -----
	\subsubsection{Section 121:7} % *DC121-7 
	% -----
		\label{DC121-7}

			\footnote{%
				% \label{}%
				This verse starts the second chunk of the section, which is verses 7--25. In the original that's \hyperlink{EMS-1839-JS-20-MAR-1839-o}{o--t}
				}%
		My son, peace be unto thy soul; thine adversity and thine afflictions shall be but a small
			\footnote{%
				% \label{}%
				See \hyperref[2CORINTHIANS-4-17]{2 Corinthians 4:17}.
				}%
		moment;

		% \begin{framed}
		% \scriptsize
			% \begin{compactdesc}
				% \item[JSP] 
				% % \item[BC] 
				% % \item[]
			% \end{compactdesc}
		% \end{framed}

	% -----
	\subsubsection{Section 121:8} % *DC121-8 
	% -----
		\label{DC121-8}

		And then, if thou endure it well, God shall exalt thee on high; thou shalt triumph over all thy foes.

		% \begin{framed}
		% \scriptsize
			% \begin{compactdesc}
				% \item[JSP] 
				% % \item[BC] 
				% % \item[]
			% \end{compactdesc}
		% \end{framed}

	% -----
	\subsubsection{Section 121:9} % *DC121-9 
	% -----
		\label{DC121-9}

		Thy friends do stand by thee, and they shall hail thee again with warm hearts and friendly hands.

		% \begin{framed}
		% \scriptsize
			% \begin{compactdesc}
				% \item[JSP] 
				% % \item[BC] 
				% % \item[]
			% \end{compactdesc}
		% \end{framed}

	% -----
	\subsubsection{Section 121:10} % *DC121-10 
	% -----
		\label{DC121-10}

		Thou art not yet as Job; thy friends do not contend against thee, neither charge thee with transgression, as they did Job.

		% \begin{framed}
		% \scriptsize
			% \begin{compactdesc}
				% \item[JSP] 
				% % \item[BC] 
				% % \item[]
			% \end{compactdesc}
		% \end{framed}

	% -----
	\subsubsection{Section 121:11} % *DC121-11 
	% -----
		\label{DC121-11}

		And they who do charge thee with transgression, their hope shall be blasted, and their prospects shall melt away as the hoar frost melteth%
			\footnote{%
				% \label{}%
				See \hyperref[EXODUS-16-14]{Exodus 16:14}; there is a theme here in this section that relates back to Exodus 16, since it speaks there of the children of Israel gathering the manna from the ground after the frost disappears. As it says in \hyperref[DC121-20]{verse 20}, the people who are not righteous will not be able to gather anything after the frost disappears. This parallel deserves to be explored further, both in this section and throughout this letter.
				}
		before the burning rays of the rising sun;

		% \begin{framed}
		% \scriptsize
			% \begin{compactdesc}
				% \item[JSP] 
				% % \item[BC] 
				% % \item[]
			% \end{compactdesc}
		% \end{framed}

	% -----
	\subsubsection{Section 121:12} % *DC121-12 
	% -----
		\label{DC121-12}

		And also that God hath set his hand and seal to change the times and seasons, and to blind their minds, that they may not understand his marvelous workings; that he may prove them also and take them in their own craftiness;

		% \begin{framed}
		% \scriptsize
			% \begin{compactdesc}
				% \item[JSP] 
				% % \item[BC] 
				% % \item[]
			% \end{compactdesc}
		% \end{framed}

	% -----
	\subsubsection{Section 121:13} % *DC121-13 
	% -----
		\label{DC121-13}

		Also because their hearts are corrupted, and the things which they are willing to bring upon others, and love to have others suffer, may come upon themselves to the very uttermost;

		% \begin{framed}
		% \scriptsize
			% \begin{compactdesc}
				% \item[JSP] 
				% % \item[BC] 
				% % \item[]
			% \end{compactdesc}
		% \end{framed}

	% -----
	\subsubsection{Section 121:14} % *DC121-14 
	% -----
		\label{DC121-14}

		That they may be disappointed also, and their hopes may be cut off;

		% \begin{framed}
		% \scriptsize
			% \begin{compactdesc}
				% \item[JSP] 
				% % \item[BC] 
				% % \item[]
			% \end{compactdesc}
		% \end{framed}

	% -----
	\subsubsection{Section 121:15} % *DC121-15 
	% -----
		\label{DC121-15}

		And not many years hence, that they and their posterity shall be swept from under heaven, saith God, that not one of them is left to stand by the wall.

		% \begin{framed}
		% \scriptsize
			% \begin{compactdesc}
				% \item[JSP] 
				% % \item[BC] 
				% % \item[]
			% \end{compactdesc}
		% \end{framed}

	% -----
	\subsubsection{Section 121:16} % *DC121-16 
	% -----
		\label{DC121-16}

		Cursed are all those that shall lift up the heel against mine anointed, saith the Lord, and cry they have sinned when they have not sinned before me, saith the Lord, but have done that which was meet in mine eyes, and which I commanded them.

		% \begin{framed}
		% \scriptsize
			% \begin{compactdesc}
				% \item[JSP] 
				% % \item[BC] 
				% % \item[]
			% \end{compactdesc}
		% \end{framed}

	% -----
	\subsubsection{Section 121:17} % *DC121-17 
	% -----
		\label{DC121-17}

		But those who cry transgression do it because they are the servants of sin, and are the children of disobedience themselves.

		% \begin{framed}
		% \scriptsize
			% \begin{compactdesc}
				% \item[JSP] 
				% % \item[BC] 
				% % \item[]
			% \end{compactdesc}
		% \end{framed}

	% -----
	\subsubsection{Section 121:18} % *DC121-18 
	% -----
		\label{DC121-18}

		And those who swear falsely against my servants, that they might bring them into bondage and death---

		% \begin{framed}
		% \scriptsize
			% \begin{compactdesc}
				% \item[JSP] 
				% % \item[BC] 
				% % \item[]
			% \end{compactdesc}
		% \end{framed}

	% -----
	\subsubsection{Section 121:19} % *DC121-19 
	% -----
		\label{DC121-19}

		Wo unto them; because they have offended my little ones they shall be severed from the ordinances of mine house.

		% \begin{framed}
		% \scriptsize
			% \begin{compactdesc}
				% \item[JSP] 
				% % \item[BC] 
				% % \item[]
			% \end{compactdesc}
		% \end{framed}

	% -----
	\subsubsection{Section 121:20} % *DC121-20 
	% -----
		\label{DC121-20}

		Their basket shall not be full, their houses and their barns shall perish, and they themselves shall be despised by those that flattered them.

		% \begin{framed}
		% \scriptsize
			% \begin{compactdesc}
				% \item[JSP] 
				% % \item[BC] 
				% % \item[]
			% \end{compactdesc}
		% \end{framed}

	% -----
	\subsubsection{Section 121:21} % *DC121-21 
	% -----
		\label{DC121-21}

		They shall not have right to the priesthood, nor their posterity after them from generation to generation.

		% \begin{framed}
		% \scriptsize
			% \begin{compactdesc}
				% \item[JSP] 
				% % \item[BC] 
				% % \item[]
			% \end{compactdesc}
		% \end{framed}

	% -----
	\subsubsection{Section 121:22} % *DC121-22 
	% -----
		\label{DC121-22}

		It had been better for them that a millstone had been hanged about their necks, and they drowned in the depth of the sea.

		% \begin{framed}
		% \scriptsize
			% \begin{compactdesc}
				% \item[JSP] 
				% % \item[BC] 
				% % \item[]
			% \end{compactdesc}
		% \end{framed}

	% -----
	\subsubsection{Section 121:23} % *DC121-23 
	% -----
		\label{DC121-23}

		Wo unto all those that discomfort my people, and drive, and murder, and testify against them, saith the Lord of Hosts; a generation of vipers shall not escape the damnation of hell.

		% \begin{framed}
		% \scriptsize
			% \begin{compactdesc}
				% \item[JSP] 
				% % \item[BC] 
				% % \item[]
			% \end{compactdesc}
		% \end{framed}

	% -----
	\subsubsection{Section 121:24} % *DC121-24 
	% -----
		\label{DC121-24}

		Behold, mine eyes see and know all their works, and I have in reserve a swift judgment in the season thereof, for them all;

		% \begin{framed}
		% \scriptsize
			% \begin{compactdesc}
				% \item[JSP] 
				% % \item[BC] 
				% % \item[]
			% \end{compactdesc}
		% \end{framed}

	% -----
	\subsubsection{Section 121:25} % *DC121-25 
	% -----
		\label{DC121-25}

		For there is a time appointed for every man, according as his works shall be.

		% \begin{framed}
		% \scriptsize
			% \begin{compactdesc}
				% \item[JSP] 
				% % \item[BC] 
				% % \item[]
			% \end{compactdesc}
		% \end{framed}

	% -----
	\subsubsection{Section 121:26} % *DC121-26 
	% -----
		\label{DC121-26}

			\footnote{%
				% \label{}%
				This verse starts the third chunk taken from the original letter. This third chunk comprises verses 26--32. In the original that's \hyperlink{EMS-1839-JS-20-MAR-1839-bb}{bb--dd}.
				}% 
		God%
			\footnote{%
				% \label{}%
				Note that starting this chunk of material in this way is very misleading---it cuts off a part of the statement which helps understand the verse as it stands now. In the original it reads, ``We beseath of you bretheren, that \sout{bare} <\textbf{you bear}> with those [w]ho do not feel themselves more worthey than yourselves, while we Exort one another, to a reffermation [reformation?], with one an all. both old and young. teachers and taugt both high and low rich and poor---bond and free. Male and female. let honesty and sobriety, and cander and solemnity, and virtue, and pureness, and Meekness, and simplisity, Crown our heads in every place, and in fine becom as little Children without mallice guile or \sout{high packrichy} Hypokrisy: and now Bretheren after your tribulations if you do this--- things, and exercise fervent prayer, and faith in the sight of God Always, he shall give unto you knowledge by his holy spirit yea by the unspeakable gift of the holy-Ghost that has not been revealed since the world was untill now...''
				}
		shall give unto you knowledge by his Holy Spirit, yea, by the \hyperref[UNSPEAKABLE]{unspeakable} gift of the Holy Ghost, that has not been revealed since the world was until now;

		\begin{framed}
		\scriptsize
			\begin{compactdesc}
				\item[JSP] he shall give unto you knowledge by his holy spirit yea by the unspeakable gift of the holy-Ghost that has not been revealed since the world was untill now
				% \item[BC] 
				% \item[]
			\end{compactdesc}
		\end{framed}

	% -----
	\subsubsection{Section 121:27} % *DC121-27 
	% -----
		\label{DC121-27}

		Which our forefathers have awaited with anxious expectation to be revealed in the last times, which their minds were pointed to by the angels, as held in reserve for the fulness of their glory;

		\begin{framed}
		\scriptsize
			\begin{compactdesc}
				\item[JSP] which our fathers have wated with anxious expectation to be revealed in the last times which their minds were pointed to by the Angels as held in reserve for the fullness of their glory
				% \item[BC] 
				% \item[]
			\end{compactdesc}
		\end{framed}

	% -----
	\subsubsection{Section 121:28} % *DC121-28 
	% -----
		\label{DC121-28}

		A time to come in the which nothing shall be withheld, whether there be one God or many gods, they shall be manifest.

		\begin{framed}
		\scriptsize
			\begin{compactdesc}
				\item[JSP] a time to come in the which nothing shall be with held whither there be one god or many god's they shall be manifest
				% \item[BC] 
				% \item[]
			\end{compactdesc}
		\end{framed}

	% -----
	\subsubsection{Section 121:29} % *DC121-29 
	% -----
		\label{DC121-29}

		All thrones and dominions, principalities and powers, shall be revealed and set forth upon all who have endured valiantly for the gospel of Jesus Christ.

		\begin{framed}
		\scriptsize
			\begin{compactdesc}
				\item[JSP] all thrones and dominions principalities and powers shall be revealed and set forth upon all who have indured valiently for the gospel of Jesus Christ
				% \item[BC] 
				% \item[]
			\end{compactdesc}
		\end{framed}

	% -----
	\subsubsection{Section 121:30} % *DC121-30 
	% -----
		\label{DC121-30}

		And also, if there be bounds set to the heavens or to the seas, or to the dry land, or to the sun, moon, or stars---

		\begin{framed}
		\scriptsize
			\begin{compactdesc}
				\item[JSP] and also if there be bounds set to the heavens or to the seas or to the dry land or to the sun moon or starrs
				% \item[BC] 
				% \item[]
			\end{compactdesc}
		\end{framed}

	% -----
	\subsubsection{Section 121:31} % *DC121-31 
	% -----
		\label{DC121-31}

		All the times of their revolutions, all the appointed days, months, and years, and all the days of their days,%
			\footnote{%
				% \label{}%
				Saying ``days of their days'' is unusual; it reminds me of the problem in Genesis 1 of calling something a ``day'' even though the means for measuring a day hadn't been created yet.
				}
		months, and years, and all their glories, laws, and set times,%
			\footnote{%
				% \label{}%
				For this verse and the previous, see \hyperref[DC88-36]{DC 88:36--45}.
				}
		shall be revealed in the days of the dispensation of the fulness of times---

		\begin{framed}
		\scriptsize
			\begin{compactdesc}
				\item[JSP] all the times of their revolutions all their appointed days month and years and all the Days of their days, months and years and all their glories laws and set times shall be reveald in the days of the dispensation of the fullness of times
				% \item[BC] 
				% \item[]
			\end{compactdesc}
		\end{framed}

	% -----
	\subsubsection{Section 121:32} % *DC121-32 
	% -----
		\label{DC121-32}

		According to that which was ordained in the midst of the Council of the Eternal God of all other gods before this world was, that should be reserved unto the finishing and the end thereof, when every man shall enter into his eternal presence and into his immortal rest.

		\begin{framed}
		\scriptsize
			\begin{compactdesc}
				\item[JSP] according to that which was ordaind in the midst of the councyl of the eternal God of all other Gods before this world was that should be reserved unto the finishing and the end thereoff \sout{where} <when> evry man shall enter into his eternal presants and into his imortal rest
				% \item[BC] 
				% \item[]
			\end{compactdesc}
		\end{framed}

	% -----
	\subsubsection{Section 121:33} % *DC121-33 
	% -----
		\label{DC121-33}

			\footnote{%
				% \label{}%
				This verse by itself is the fourth chunk from the original source. In the original it's \hyperlink{EMS-1839-JS-20-MAR-1839-ee}{ee}.
				}%
		How long can rolling waters remain impure?  What power shall stay the heavens?  As well might man stretch forth his puny arm to stop the Missouri river in its decreed course, or to turn it up stream, as to hinder the Almighty from pouring down knowledge from heaven upon the heads of the Latter-day Saints.

		% \begin{framed}
		% \scriptsize
			% \begin{compactdesc}
				% \item[JSP] 
				% % \item[BC] 
				% % \item[]
			% \end{compactdesc}
		% \end{framed}

	% -----
	\subsubsection{Section 121:34} % *DC121-34 
	% -----
		\label{DC121-34}

			\footnote{%
				% \label{}%
				This verse starts the fifth and final chunk of the section, but note that we're now drawing from the second letter, wihch is a continuation of the first. The chunk is verses 34--46 and in the original source it's \hyperlink{EMS-1839-JS-CIR-22-MAR-1839-d}{d--h}.
				}%
		Behold, there are many called, but few are chosen.%
			\footnote{%
				% \label{}%
				See \hyperref[MATTHEW-22-8]{Matthew 22:8} (note specfically the word for ``bidden,'' \gr{κεκλημένοι}, ``those who'd been called''), \hyperref[DC95-5]{DC 95:5--6}.
				} 
		And why are they not chosen?

		\begin{framed}
		\scriptsize
			\begin{compactdesc}
				\item[JSP] Behold there are ma[n]y called but few are chosen. And why are they not chosen?
				% \item[BC] 
				% \item[]
			\end{compactdesc}
		\end{framed}

	% -----
	\subsubsection{Section 121:35} % *DC121-35 
	% -----
		\label{DC121-35}

		Because their hearts are set so much upon the things of this world, and aspire to the honors%
			\footnote{%
				% \label{}%
				Compare to the use of ``honor'' in \hyperref[MOSES-4-1]{Moses 4:1}; see notes there as well.
				}
		of men, that they do not learn this one lesson---

		\begin{framed}
		\scriptsize
			\begin{compactdesc}
				\item[JSP] Because their hearts are set so much upon the things of this world and aspire to the honors of men that they do not learn this one lesson
				% \item[BC] 
				% \item[]
			\end{compactdesc}
		\end{framed}

	% -----
	\subsubsection{Section 121:36} % *DC121-36
	% -----
		\label{DC121-36}

		That the rights of the priesthood are inseparably connected with the powers of heaven,%
			\footnote{%
				% \label{}%
				Compare \hyperref[MARK-13-25]{Mark 13:25}.
				}
		and that the powers of heaven cannot be controlled nor handled only upon the principles of righteousness.
		
		\begin{framed}
		\scriptsize
			\begin{compactdesc}
				\item[JSP] that the rights of priesthood are inseperably connected with the powers of heaven and that the powers of heaven connot be controled nor handled only upon the principals of rightiousness.
				% \item[BC] 
				% \item[]
			\end{compactdesc}
		\end{framed}
		
	% -----
	\subsubsection{Section 121:37} % *DC121-37
	% -----
		\label{DC121-37}

		That they may be conferred upon us, it is true; but when we undertake to cover our sins, or to gratify our pride, our vain ambition, or to exercise control or dominion or compulsion upon the souls of the children of men, in any degree of unrighteousness, behold, the heavens withdraw themselves; the Spirit of the Lord is grieved; and when it is withdrawn, Amen to the priesthood or the authority of that man.
		
		\begin{framed}
		\scriptsize
			\begin{compactdesc}
				\item[JSP] That they may be confered upon us it is true but when we undertake to cover our sins or to gratify our pride or vain ambition or to exercise controle or dominion or compulsion upon the souls of the children of men in any degree of unritiousness behold the heavens withdraw themselves the spirit of the Lord is grieved and when it has withdrawn \uline{Amen} to the \uline{priesthood} or the authority of that man
				% \item[BC] 
				% \item[]
			\end{compactdesc}
		\end{framed}
		
	% -----
	\subsubsection{Section 121:38} % *DC121-38
	% -----
		\label{DC121-38}
		
		Behold, ere he is aware, he is left unto himself, to kick against the pricks,%
			\footnote{%
				% \label{}%
				See \hyperref[ACTS-9-5]{Acts 9:5} and \hyperref[ACTS-26-14]{26:14}.
				}
		to persecute the saints, and to fight against God.
		
		\begin{framed}
		\scriptsize
			\begin{compactdesc}
				\item[JSP] behold ere he is aware he is left unto himself to kicken against the pricks to persecute the saints and to fight against God.
				% \item[BC] 
				% \item[]
			\end{compactdesc}
		\end{framed}
		
	% -----
	\subsubsection{Section 121:39} % *DC121-39
	% -----
		\label{DC121-39}
		
		We have learned by sad experience that it is the nature and disposition of almost all men, as soon as they get a little authority, as they suppose, they will immediately begin to exercise unrighteous dominion.
		
		\begin{framed}
		\scriptsize
			\begin{compactdesc}
				\item[JSP] We have learned by sad experiance that it is the nature and disposition of almost all men as soon as they get a little authority as they suppose they will imediately begin to [e]xercise unritious dominion.
				% \item[BC] 
				% \item[]
			\end{compactdesc}
		\end{framed}
		
	% -----
	\subsubsection{Section 121:40} % *DC121-40
	% -----
		\label{DC121-40}
		
		Hence many are called, but few are chosen.
		
		\begin{framed}
		\scriptsize
			\begin{compactdesc}
				\item[JSP] hence ma[n]y [are] called but few are ch[osen.]
				% \item[BC] 
				% \item[]
			\end{compactdesc}
		\end{framed}
		
	% -----
	\subsubsection{Section 121:41} % *DC121-41
	% -----
		\label{DC121-41}
		
		No power or influence can or ought to be maintained by virtue of the priesthood, only by persuasion, by long-suffering, by gentleness and meekness, and by love unfeigned;
		
		\begin{framed}
		\scriptsize
			\begin{compactdesc}
				\item[JSP] [No power or in]f[luence] can or ought to be maintained by <[vi]rt[ue]> of the Priesthood only by persuasion by long suffering by gentleness and meekness and by love unfaigned
				% \item[BC] 
				% \item[]
			\end{compactdesc}
		\end{framed}
		
	% -----
	\subsubsection{Section 121:42} % *DC121-42
	% -----
		\label{DC121-42}
		
		By kindness, and pure knowledge, which shall greatly enlarge the soul%
			\footnote{%
				% \label{}%
				A reference to \hyperref[ALMA-32-28]{Alma 32:28}.
				}
		without hypocrisy, and without guile---
		
		\begin{framed}
		\scriptsize
			\begin{compactdesc}
				\item[JSP] by kindness by pure knowledge which shall greatly enlarge the soul without hypocrisy and without guile
				% \item[BC] 
				% \item[]
			\end{compactdesc}
		\end{framed}
		
	% -----
	\subsubsection{Section 121:43} % *DC121-43
	% -----
		\label{DC121-43}
		
		Reproving \hyperref[BETIMES]{betimes} with sharpness, when moved upon by the Holy Ghost; and then showing forth afterwards an increase of love toward him whom thou hast reproved, lest he esteem thee to be his enemy;
		
		\begin{framed}
		\scriptsize
			\begin{compactdesc}
				\item[JSP] reproving be-times with sharpness when moved upon by the Holy Ghost and then showing forth afterwords an increas of love toward him whom thou hast reproved lest he esteem the[e] to be his enimy
				% \item[BC] 
				% \item[]
			\end{compactdesc}
		\end{framed}
		
	% -----
	\subsubsection{Section 121:44} % *DC121-44
	% -----
		\label{DC121-44}
		
		That he may know that thy faithfulness is stronger than the cords of death.
		
		\begin{framed}
		\scriptsize
			\begin{compactdesc}
				\item[JSP] that he may know that thy faithfulness is stronger than the cords of death
				% \item[BC] 
				% \item[]
			\end{compactdesc}
		\end{framed}
		
	% -----
	\subsubsection{Section 121:45} % *DC121-45
	% -----
		\label{DC121-45}
		
		Let thy \hyperref[BOWELS]{bowels}%
			\footnote{%
				% \label{}%
				Notice that the ``bowels'' are spoken of here as the seat of charity, in a similar way to how the construction ``bowels of mercy'' is used to talk about the seat of mercy. See meaning \#3 in the \hyperref[BOWELS-1828]{1828}, and also note the other occurrences of ``bowels'' earlier in this section.
				}
		also be full of charity towards all men, and to the household%
			\footnote{%
				% \label{}%
				See \hyperref[GALATIANS-6-10]{Galatians 6:10},
				}
		of faith, and let \hyperref[VIRTUE]{virtue}%
			\footnote{%
				% \label{}%
				Here, perhaps meaning \#3 from the 1828, ``Moral goodness; the practice of moral duties and the abstaining from vice, or a conformity of life and conversation to the moral law.'' But it could be others as well.
				}
		\hyperref[GARNISH]{garnish}%
			\footnote{%
				% \label{}%
				Compare \hyperref[MATTHEW-12-44]{Matthew 12:44} and \hyperref[LUKE-11-25]{Luke 11:25}, which are forms of \gr{κοσμέω}; I think this is actually a pretty helpful comparison.
				}
		thy thoughts unceasingly; then shall thy confidence%
			\footnote{%
				% \label{}%
				Possibly \gr{παρρησία} in Greek; see also \gr{πεποίθησις}, but probably not \gr{ὑπόστασις}. Compare \gr{ἀναίδεια} in \hyperref[LUKE-11-8]{Luke 11:8}, along with the related verb \gr{ἀναιδεύομαι} (``\textbf{to be unconcerned about convention, \textit{be unabashed, bold}}'') and adjective \gr{ἀναιδής} (``\textbf{\textit{shameless, bold}}''), but note that neither the verb nor adjective appears in the NT; they're only found elsewhere.
				}
		wax%
			\footnote{%
				% \label{}%
				Compare \hyperref[PHILIPPIANS-1-14]{Philippians 1:14}.
				}
		strong in the presence of God; and the doctrine%
			\footnote{%
				% \label{}%
				A knowledge of the laws? Notice how this doctrine, whatever it is, accumulates gradually, like the dew; the intransitive use of ``distil,'' which seems to be the one used here (unless the object of the transitive verb is suppressed), does not include the idea of purification or getting rid of impurities, as we see in the transitive use, at least according to the 1828. But meaning \#6 in the OED is clearer, and it's for the intransitive use: ``To become vaporized and then condensed into liquid; to undergo distillation; to drop, pass, or condense from the still. \textit{\textbf{to distil over}}: to pass over in the form of vapour which again condenses into a liquid.'' So the doctrine gets \textit{purified} over time---more and more focused on what is essential to the priesthood, I would say.
				}
		of the priesthood shall \hyperref[DISTIL]{distil}%
			\footnote{%
				% \label{DC121-45-distil}%
				``Distil'' occurs only three times in all the scriptures; here, at \hyperref[DEUTERONOMY-32-2]{Deuteronomy 32:2}, and \hyperref[JOB-36-28]{Job 36:28}. This passage in verse 45 seems like a clear reference to the Deuteronomy verse, which in the KJV reads: ``My doctrine shall drop as the rain, my speech shall distil as the dew, as the small rain upon the tender herb, and as the showers upon the grass:''. Check out the verse in Deuteronomy for more. I also want to connect this to the idea of receiving ``line upon line, precept on precept,'' as from \hyperref[2NEPHI-28-30]{2 Nephi 28:30} and \hyperref[DC98-12]{DC 98:12}.
				}
		upon thy soul as the dews%
			\footnote{%
				% \label{}%
				In some \hyperlink{EMS-1842-JS-16-23-AUG-1842-cc}{blessings} given in August 1842, JS uses a similar manner of speaking: ``As the dews upon Mount Hermon may the distillations of thy divine grace, glory and honor in the plenitude of thy mercy, and power and goodness be poured down upon the head of thy servant.'' He's referencing \hyperref[PSALMS-133-3]{Psalms 133:3}. See also \hyperref[DANIEL-4-15]{Daniel 4:15} and notes there; \hyperref[DC128-19]{DC 128:19} (``dews of Carmel''); 
				}
		from%
			\footnote{%
				% \label{}%
				In this verse the dews are spoken of as ``distilling'' upon ours souls ``from heaven.'' Dew doesn't actually come from heaven; it comes from condensation of water vapor in the air when the temperature drops. We can extend the metaphor in that direction of the condensation of what is already around us or what we are already breathing in or experiencing---it changes form so that we understand and experience it in a new way, even a more useful way, if you need the water for certain purposes. The ``from heaven'' bit also reminds me of manna, as spoken of in \hyperref[EXODUS-16-4]{Exodus 16:4}. I
				}
		heaven.
	
		\begin{framed}
		\scriptsize
			\begin{compactdesc}
				\item[JSP] thy bowells also being full of charity towards all men and to the household of faith and virtue garnish thy thoughts unceasingly then shall thy confidence wax strong in the presants of God, and the doctrins of the Priesthood shall destill upon thy soul as the dews from heaven
				% \item[BC] 
				% \item[]
			\end{compactdesc}
		\end{framed}
		
	% -----
	\subsubsection{Section 121:46} % *DC121-46
	% -----
		\label{DC121-46}
		
		The Holy Ghost shall be thy constant companion, and thy scepter%
			\footnote{%
				% \label{}%
				
				
				See also \hyperref[DC85-7]{DC 85:7}, 
				}
		an unchanging scepter of righteousness and truth; and thy dominion shall be an everlasting dominion, and without compulsory means it shall flow unto%
			\footnote{%
				% \label{}%
				See \hyperref[DC111-8]{DC 111:8}, which is the only place in the scriptures outside of OT-related verses that you get the phrase ``flow unto'' (``flow to'' does not appear at all). 111:8 is pretty significant in understanding this verse, I think, since it's speaking of the ``peace and power of my Spirit'' that will flow. The other places are \hyperref[ISAIAH-2-2]{Isaiah 2:2}, \hyperref[MICAH-4-1]{Micah 4:1}, and \hyperref[2NEPHI-12-2]{2 Nephi 12:2}, which comes from the Isaiah verse. ``Flow together'' appears three times: \hyperref[ISAIAH-60-5]{Isaiah 60:5}, \hyperref[JEREMIAH-31-12]{Jeremiah 31:12}, and \hyperref[JEREMIAH-51-44]{Jeremiah 51:44}. See also \hyperref[1NEPHI-20-18]{1 Nephi 20:18}, 
				}
		thee forever and ever.
		
		\begin{framed}
		\scriptsize
			\begin{compactdesc}
				\item[JSP] the Holy Ghost shall be thy constant companion and thy septer an unchanging septer of ritiousness and truth and thy dominion shall be an everlasting dominion and without compulsory means it shall flow [un]to thee for ever and ever.
				% \item[BC] 
				% \item[]
			\end{compactdesc}
		\end{framed}